\track{T4\quad 广义复几何、sigma 模型与膜}{本轨道 18 条：A/B 膜是物理侧「能装 Hecke 算子」的容器，它们翻译成数学就是相干层与导出范畴里的对象。}

\begin{word}{sigma 模型}{sigma model}{sigma-model}{§5 p.65}
\mathdef{给定光滑流形 $X$（靶空间）与二维曲面 $\Sigma$（世界面），sigma 模型的基本场是映射 $\Phi:\Sigma\to X$；经典作用量是把 $X$ 上的度量与 $2$-形式沿 $\Phi$ 拉回后积分，例如 $\int_\Sigma \Phi^* g$ 与 $\int_\Sigma\Phi^*\omega$。拓扑版本只保留一个标量超荷 $Q$，并要求作用量可写成 $\{Q,V\}$ 加一个拓扑项；于是物理量取 $Q$-上同调，$\Sigma$ 的度量与复结构都不进入答案，真正重要的是 $X$ 上被 $Q$ 看见的几何：复结构、辛结构或广义复结构。}
\mathex{取 $X=T^2$、$\Sigma=T^2$。每一张映射的同伦类由 $\pi_1(T^2)$ 两条生成元的像决定；在 A-模型中要求 $\Phi$ 全纯，写成 $z\mapsto a z+b$，其中 $a\in\mathbb Z$、$b\in T^2$，所以对每个 $a$，映射的模空间是一份 $T^2$（实维 $2$）。}
\physdef{KW06 的用法是：四维扭化 $\mathcal N=4$ 超对称杨--米尔斯理论紧化在 $\Sigma\times C$ 上（$C$ 是黎曼面），低能有效理论就是以 Hitchin 模空间 $\mathcal M_H(G,C)$ 为靶空间的二维拓扑 sigma 模型。它自动带 $(4,4)$ 超对称，因为 $\mathcal M_H$ 超凯勒；选取复结构对 $(J_+,J_-)$ 就得到不同的拓扑模型。}
\physex{取 $G=SU(2)$、$C$ 亏格 $2$，则 $\mathcal M_H$ 复维为 $2\dim_{\mathbb C}\mathrm{Bun}_G=2\cdot 3\cdot(2-1)=6$，实维 $12$。世界面取 $\Sigma=T^2$ 时，模型是 $T^2$ 上的 $(4,4)$ 拓扑场论，其 $Q$-上同调给出膜范畴的态空间。}
\bridge{数学侧给 $X$ 一个几何结构（复结构、辛结构、超凯勒三元组），物理侧就得到一个拓扑 sigma 模型；$\Phi$ 是场，$\Sigma$ 的边界条件是膜，$Q$-上同调则是导出范畴中的态空间。KW06 的全部字典都建立在这条「几何结构映到模型类型」的规则上（\S5.1）。}
\intuit{sigma 模型就是「世界面在靶空间里画图」，拓扑版本把连续形变都当成 $Q$-恰当而丢掉，只记得画图的同伦类。靠得住是因为拓扑超荷使作用量减去拓扑项后整体 $Q$-恰当。}
\pitfall{不要把它当成一般的二维量子场论或弦论：没有扭化与超对称就没有拓扑性；KW06 用的也不是任意 sigma 模型，而是靶空间超凯勒、带 $(4,4)$ 超对称的二维版本。}
\end{word}

\begin{word}{A-模型}{A-model}{a-model}{§1 p.2；§5.1 p.65}
\mathdef{设 $X$ 是辛流形，辛形式为 $\omega$，并带闭 $2$-形式 $B$（$B$-场）。$X$ 上的 A-模型是只依赖复化 Kähler 类 $[B+i\omega]$ 的拓扑 sigma 模型：把 $\omega$ 换成 $\omega+d\eta$、$B$ 换成 $B+d\lambda$ 而不改上同调类时，关联函数不变。它的数学化身是辛几何一侧：A-膜以余各向同性子流形加酉丛为对象，态空间由开弦谱（Floer 型理论）给出。}
\mathex{取 $X=T^2$、$\omega$ 为面积形式、$B=0$。$T^2$ 上 $H^2$ 是一维的，所以 A-模型的唯一连续参数是复化 Kähler 模 $\tau=i\,\mathrm{vol}$。世界面取 $\Sigma=S^1\times\mathbb R$ 时，A-膜是 $T^2$ 中的圆周 $L=S^1$（Lagrange 子流形）加上平坦 $U(1)$ 联络 $\nabla=d+i\theta\,dt$，$\theta\in\mathbb R/\mathbb Z$。}
\physdef{在 KW06 的二维图像里，A-模型对应 $J_+=-J_-=\hat J$：未破缺方程 $(1-i\hat J)\bar\partial\Phi=0$ 与 $(1+i\hat J)\partial\Phi=0$ 互为共轭，等价于 $\bar\partial\Phi=0$，即 $\Phi$ 在 $\hat J$ 中全纯。原文特意指出：它实际依赖的是与 $\hat J$ 相联的辛形式 $\omega_{\hat J}$，而不是复结构本身（\S5.1）。从四维看，它是 $R$-对称扭化后的四维理论经紧化、取 $Q$-上同调得到的二维模型，所以 A/B 模型是「四维理论的二维影子」，靶空间由四维理论的模空间给出。}
\physex{取四维扭化参数 $t=1$（实数），复结构对为 $(w_+,w_-)=(-1,1)$，模型是复结构 $K$ 中的 A-模型；$t=0$ 时得到复结构 $I$ 中的 A-模型。$G=SU(2)$、$C$ 亏格 $2$ 时靶空间 $\mathcal M_H$ 复维 $6$，可取辛形式 $\omega_K$。}
\bridge{A-模型对应辛结构，B-模型对应复结构。$S$-对偶作用在典范参数上为 $\Psi\mapsto-1/\Psi$，交换两者；在二维它表现为复结构 $J$ 中的 B-模型与复结构 $K$ 中的 A-模型互换（\S5.2--\S5.3 与 Table~2）。}
\intuit{A 是「面积」：A-模型只关心映射扫过靶空间时的辛面积。靠得住是因为作用量中只剩 $\int_\Sigma\Phi^*\omega$ 这类拓扑项是 $Q$-闭而非恰当的。}
\pitfall{A-模型不是 A-膜，也不是「在某个复结构中的模型」的字面意思：前者是世界面上的二维场论，后者是对靶空间的边界条件；原文说「复结构 $\hat J$ 中的 A-模型」只是简称，它依赖的是 $\omega_{\hat J}$。}
\end{word}

\begin{word}{B-模型}{B-model}{b-model}{§1 p.4；§5.1 p.65}
\mathdef{设 $X$ 是复流形，复结构为 $J$。$X$ 上的 B-模型是基本场 $\Phi:\Sigma\to X$、场方程为 $\bar\partial\Phi=0$（$\Phi$ 全纯）的拓扑 sigma 模型；作用量只依赖复结构，与 Kähler 类无关，原因是超荷 $\bar\partial$ 只用 $J$ 定义，而作用量中的 Kähler 部分在扭化后是 $Q$-恰当的，度量与 Kähler 模因此不进入关联函数。它的数学化身是复几何一侧：B-膜的范畴是 $X$ 上相干层（coherent sheaves）的导出范畴 $D^b(\mathrm{Coh}(X))$，对象是层或层复形，态空间是 $\mathrm{Ext}^\bullet$，可用 Dolbeault 复形计算。}
\mathex{取 $X=\mathbb P^1$。线丛 $\mathcal O(k)$ 给出一族 B-膜；点层 $\mathcal O_p$ 给出支撑在点上的 B-膜。$\mathrm{Ext}^\bullet(\mathcal O_p,\mathcal O_p)=\wedge^\bullet T_p\mathbb P^1$ 是一维的，其中 $\mathrm{Ext}^1(\mathcal O_p,\mathcal O_p)=\mathbb C$ 正对应「点可以沿 $\mathbb P^1$ 移动」这个一维形变；又由 Serre 对偶，$\mathrm{Ext}^1(\mathcal O,\mathcal O(-2))=H^1(\mathbb P^1,\mathcal O(-2))=\mathbb C$。取 $X=T^2$、$\tau=i$：B-膜由全纯线丛与点层给出，线丛对应填满空间的膜，点层对应零膜。}
\physdef{$J_+=J_-=\hat J$ 时要求 $w_+=w_-$，即 $t=\pm i$，所得模型是复结构 $\pm J$ 中的 B-模型（\S5.1）。B-模型对 $\mathrm{Im}\,\tau$ 与控制 Kähler 类的参数不敏感，但 $B$-场的 $(0,2)$ 分量会改变 B-膜的范畴（把相干层扭化）。}
\physex{取 $t=i$ 得 $\Psi=\infty$，模型是复结构 $J$ 中的 B-模型；$J$ 与 $\Omega_J$ 都是平坦 $G_{\mathbb C}$ 丛的复结构，不依赖 $C$ 的复结构，所以这族模型与 $C$ 的模无关。$G=SU(2)$、$C$ 亏格 $2$ 时，它是 $\mathrm{Hom}(\pi_1(C),G_{\mathbb C})/G_{\mathbb C}$ 上的 sigma 模型。}
\bridge{对象级：层 $\mathcal F$ 映到支撑在 $\mathrm{supp}\,\mathcal F$ 上、Chan-Paton 丛为 $\mathcal F$ 的 B-膜。态射级：$\mathrm{Hom}(\mathcal F,\mathcal G)=\mathrm{Ext}^\bullet_{\mathcal O_X}(\mathcal F,\mathcal G)$ 就是两端分别落在两个膜上的开弦态空间，分次由鬼数给出，复合对应开弦算子的乘积展开。Fourier--Mukai 变换以核层实现，是 $T$-对偶的数学化（\S5.1、\S9）。}
\intuit{B 是「全纯」：B-模型只问映射是不是全纯的。靠得住是因为作用量只含复结构，$\bar\partial$ 方程根本不涉及度量。}
\pitfall{B-模型不是 B-膜，也不是「复流形」本身；说它「只依赖复结构」不等于与 $B$-场无关：$B$ 的 $(0,2)$ 分量会扭化 Chan-Paton 层，从而改变膜的范畴。}
\end{word}

\begin{word}{镜像对称}{mirror symmetry}{mirror-symmetry}{§5.3 p.72；§5.4 p.74}
\mathdef{镜像对称最干净的数学陈述是一对复流形 $X$ 与 $\hat X$（通常为 Calabi--Yau）之间 A-模型与 B-模型的等价：$X$ 的 A-模型等价于 $\hat X$ 的 B-模型，且 $X$ 的复结构模与 $\hat X$ 的复化 Kähler 模互换。范畴形式（同调镜像对称）断言 $D^b(\mathrm{Coh}(\hat X))$ 与 $X$ 的 A-膜范畴（Fukaya 型）等价；Hodge 数按 $h^{p,q}(X)=h^{d-p,q}(\hat X)$ 互换。}
\mathex{椭圆曲线给出最简单的例子：$E_\tau$ 与 $E_{-1/\tau}$ 互为镜像，复结构模 $\tau$ 与复化 Kähler 模互换。$T^2$ 上的镜像对称可以逐点写成 $\tau\leftrightarrow\rho$（$\rho$ 是复化 Kähler 参数）；更高维的例子是五次三维 Calabi--Yau 与它的 Batyrev 镜像。}
\physdef{物理上镜像对称就是「$X$ 的 A-模型等于 $\hat X$ 的 B-模型」这条等价。KW06 的入口是 $S$-对偶：四维 $S$-对偶在二维表现为镜像对称，把 $\mathcal M_H(G,C)$ 与 $\mathcal M_H({}^LG,C)$ 配成镜像对（\S5.3--\S5.4）。}
\physex{取 $G=SU(2)$、${}^LG=SO(3)$、$C$ 亏格 $2$：在复结构 $J$ 中，$\mathcal M_H(SU(2))$ 与 $\mathcal M_H(SO(3))$ 的 B-模型互为镜像；两个 Hitchin 纤维化同底 $B$，对应纤维是互为对偶的复环面（谱曲线的广义 Prym 簇）。}
\bridge{镜像对称等于在 Hitchin 纤维的环面上作 $T$-对偶、同时在基上加一个映射 $\Xi$；这正是 SYZ 图景，也是 Fourier--Mukai 变换的物理来源。具体操作是：$S$-对偶先在 $J$ 中给 B-模型与 $K$ 中 A-模型的等价，再复合 $\varphi\mapsto i\varphi$ 得到 $J$ 中的镜像（\S5.3）。}
\intuit{镜像对称不是「两个流形一样」，而是把一个流形的复几何换成另一个流形的辛几何，就像傅里叶变换把位置换成动量。}
\pitfall{不是几何等同：镜像流形一般不拓扑同构，Hodge 数会按 $h^{p,q}\leftrightarrow h^{d-p,q}$ 互换。KW06 的镜像对是超凯勒的 Hitchin 模空间（非紧），不是通常的紧 Calabi--Yau；SYZ 图景在原文只作启发，不是严格证明。}
\end{word}

\begin{word}{广义复几何}{generalized complex geometry}{generalized-complex-geometry}{§5.2 p.65}
\mathdef{在光滑流形 $X$ 上取 $T\oplus T^*$，配自然的不定度量 $\langle v+\xi,v'+\xi'\rangle=\frac12(\xi(v')+\xi'(v))$；$TX$ 与 $T^*X$ 都是迷向子丛、彼此配对。广义复几何就是这套结构上的复几何：一个保度量、平方为 $-1$ 且对 Courant 括号可积的自同态 $\mathcal I:T\oplus T^*\to T\oplus T^*$ 是广义复结构；理论由 Hitchin 提出、Gualtieri 系统化，目的是让复几何与辛几何成为同一理论的两个极端。}
\mathex{$X$ 带复结构 $J$ 时，$\mathcal I_J=\mathrm{diag}(J,-J^t)$ 是广义复结构；$X$ 带辛结构 $\omega$（$d\omega=0$）时，$\mathcal I_\omega=\begin{pmatrix}0&-\omega^{-1}\\\omega&0\end{pmatrix}$ 也是。取 $X=\mathbb C$、$J=i$ 与 $\omega=dx\wedge dy$，二者给出两个不同的广义复结构；在 $T^2$ 上它们分别对应复结构模与 Kähler 模，正是 $T$-对偶要交换的两边。}
\physdef{KW06 用广义复几何统一两个拓扑扭化：$(2,2)$ 超对称条件等价于存在一对相容的广义复结构，扭化后两个拓扑场论各由其中一个决定（\S5.2，原文注明这一点一般认为成立，未给完整证明）。所以「A-模型与 B-模型」不是两个互不相干的理论，而是同一族广义复结构的两个极限。}
\physex{在超凯勒的 $\mathcal M_H$ 上取 $(w_+,w_-)=(-t,t^{-1})$ 得整族广义复结构 $\mathcal I_t$（式 5.30），它只依赖 $J$ 与 $\Omega_J=\omega_K+i\omega_I$，与 $C$ 的复结构无关。$t=\pm i$ 时拓扑模型是复结构 $J$ 中的 B-模型；$t$ 取实数时是 A-模型；$t\ne\pm i$ 时可用 $\mathcal I_t=\mathcal M(B_0)\mathcal I_{\omega_0}\mathcal M(B_0)^{-1}$（式 5.31）约化为某个辛结构的 A-模型。}
\bridge{广义复几何是 A/B 模型共居的舞台：$\mathcal I_J$ 给 B-模型，$\mathcal I_\omega$ 给 A-模型，$B$-场在两者之间作规范变换 $\mathcal I\mapsto\mathcal M(B_0)\mathcal I\mathcal M(B_0)^{-1}$；范畴层面 B-膜对应相干层、A-膜对应 Fukaya 型对象（\S5.2、\S11.1）。}
\intuit{就像给流形同时装上「切向量」与「$1$-形式」两套坐标并允许它们互相旋转；复结构与辛结构只是两个极端，中间还有一整族既不全纯也不辛的对象。}
\pitfall{不是复几何与辛几何的并集：一般 $\mathcal I$ 既不是复结构也不是辛结构，理论的价值恰恰在中间那族结构。也不是只要 $\mathcal I^2=-1$ 就够，还需保度量与 Courant 可积；此外 B-场变换 $\mathcal M(B_0)$ 只在 $B_0$ 闭时保持可积性。}
\end{word}

\begin{word}{广义复结构}{generalized complex structure}{generalized-complex-structure}{§5.2 p.66}
\mathdef{$T\oplus T^*$ 上的自同态 $\mathcal I$，满足 $\mathcal I^2=-1$、关于自然配对保度量，并且可积：$\mathcal I$ 的 $+i$ 特征子丛 $L\subset(T\oplus T^*)\otimes\mathbb C$ 是极大迷向的，且对 Courant 括号封闭。等价地，若记 $\mathcal I=\begin{pmatrix}A&\pi\\\sigma&-A^t\end{pmatrix}$，则可积性给出一组关于 $A$、$\pi$、$\sigma$ 的微分方程；$\mathcal I_J$ 与 $\mathcal I_\omega$ 分别对应「$\pi=0$」与「$A=0$、$\sigma=\omega$ 闭」两种极端。}
\mathex{$X$ 为复流形：$\mathcal I_J=\mathrm{diag}(J,-J^t)$ 的 $+i$ 特征空间是 $T^{0,1}\oplus (T^*)^{1,0}$ 的复共轭，可积性就是 $J$ 的 Nijenhuis 张量为零。$X$ 为辛流形：$\mathcal I_\omega=\begin{pmatrix}0&-\omega^{-1}\\\omega&0\end{pmatrix}$，可积性就是 $d\omega=0$；其 $+i$ 特征子空间由 $v-i\omega(v,\cdot)$ 张成。取 $B$ 闭，$\mathcal M(B)=\begin{pmatrix}1&0\\B&1\end{pmatrix}$ 与 $\mathcal I$ 共轭仍得广义复结构。}
\physdef{在 KW06 里，一个拓扑模型由一个广义复结构决定：B-模型由 $\mathcal I_J$、A-模型由 $\mathcal I_\omega$ 决定；两个相容广义复结构的存在性就是 $(2,2)$ 超对称（\S5.2）。$B$-场开启时用 $\mathcal M(B_0)$ 作用，得到「带 $B$-场的 $A$-模型」对应的广义复结构（Gualtieri 命题 4.34）。}
\physex{$\mathcal M_H$ 上的一族 $\mathcal I_t$（式 5.30）给出从复结构型到辛结构型的连续插值：$t=\pm i$ 是纯复的 $\mathcal I_J$，$t$ 实数时是辛型的 $\mathcal I_{\omega}$（$t=\pm1$ 对应 $\omega_K$，$t=0,\infty$ 对应 $\omega_I$）。它只依赖 $J$、$\omega_I$、$\omega_K$，因而与 $C$ 的模无关；这是几何 Langlands 中「把 $C$ 的依赖全放进基 $B$」的技术基础。}
\bridge{广义复结构是 A-模型与 B-模型的「共同母体」：$\mathcal I_J$ 与 $\mathcal I_\omega$ 分别对应复结构与辛结构，S-对偶与 T-对偶则作用在这族结构上，把一种类型换成另一种。数学上它等价于一个复 Lie 代数胚（Dirac 结构）的极大迷向子丛，因此也是 Poisson 几何、Dirac 结构的统一语言（\S5.2）。}
\intuit{把复结构想成「把切空间分成 $(1,0)$ 与 $(0,1)$」；广义复结构则把切空间与余切空间一起分成两半，于是全纯切向量与全纯 $1$-形式可以互相转化。}
\pitfall{$\mathcal I_J$ 与 $\mathcal I_\omega$ 只是两个极端例子，不是全部；一般 $\mathcal I$ 既非复也非辛。另外 $\mathcal I^2=-1$ 只给近广义复结构，可积性（$+i$ 子丛对 Courant 括号封闭）是独立条件，正如近复结构不等于复结构。}
\end{word}

\begin{word}{B-场}{B-field}{b-field}{§4.1 p.46}
\mathdef{靶空间 $X$ 上的闭 $2$-形式 $B$（整体上更准确的说法是 gerbe 的联络，局部才是 $2$-形式），与辛形式 $\omega$ 并列进入二维作用量：$\int_\Sigma\Phi^*(\omega-iB)$。物理只依赖复化 Kähler 类 $[B+i\omega]\in H^2(X,\mathbb C)/H^2(X,\mathbb Z)$；把 $B$ 换成 $B+d\lambda$ 是规范变换，不改变任何物理量。它在数学上的两种作用是：平移复化 Kähler 模，以及通过 $\mathcal M(B_0)$ 对广义复结构作 B-变换。}
\mathex{取 $X=T^2$、$\omega=\mathrm{vol}$、$B=b\,\mathrm{vol}$，$b\in\mathbb R/\mathbb Z$：复化 Kähler 模是 $\rho=b+i\,\mathrm{vol}$，大规范变换 $b\mapsto b+1$ 保持物理不变。$X=T^*M$ 上取 $B=\pi^*\beta$（$\beta$ 是 $M$ 上的闭 $2$-形式），则 $B$ 沿纤维为零，是一类典型的闭 $B$-场。广义复结构的 B-变换是 $\mathcal M(B_0)=\begin{pmatrix}1&0\\B_0&1\end{pmatrix}$，$\mathcal I(B_0)=\mathcal M(B_0)\mathcal I\mathcal M(B_0)^{-1}$。}
\physdef{在 KW06 的四维理论里，$B$-场来自 $\theta$ 角与 $\tau$：$\theta$ 角给 $B'=-\frac{\theta}{2\pi}\omega_I=-(\mathrm{Re}\,\tau)\omega_I$；把 $\omega$ 与 $B$ 的组合称为复化 Kähler 类，A-模型只依赖 $[B+i\omega]$。四维 $S$-对偶把 $\tau\mapsto-1/\tau$，同时把 $B+i\omega$ 送到对偶流形上（\S5.1--\S5.3）。}
\physex{取 $\mathcal M_H$ 与辛形式 $\omega=(\mathrm{Im}\,\tau)\omega_K$，则 $B=-(\mathrm{Re}\,\tau)\omega_I$ 与 $\omega$ 一起进入 A-模型；A-模型真正依赖的参数是 $\Psi=\mathrm{Re}\,\tau+i\,\mathrm{Im}\,\tau\frac{t-t^{-1}}{t+t^{-1}}$（式 5.34）。在膜上情况更精细：支撑在 Lagrange 子流形 $N$ 上的 Chan-Paton 线丛 $L$ 给 A-膜，当且仅当 $F+B|_N=0$，这使 A-膜对 $B$ 的所有 Hodge 分量都敏感（\S5.1）。}
\bridge{$B$-场在数学侧有三个化身：复化 Kähler 类 $[B+i\omega]$（A-模型的参数）、广义复结构的 B-变换 $\mathcal M(B_0)$、以及平坦 $B$-场 $b\in H^2(X,U(1))$ 把 Chan-Paton 丛扭化成扭化层（§7、§11.1）。在 B-膜一侧，$B$ 的 $(0,2)$ 分量改变相干层的范畴；在 A-膜一侧，$B$ 与曲率 $F$ 只以 $F+B$ 出现。}
\intuit{$B$-场像穿过靶空间的磁场通量：弦的世界面只与它的面积分耦合，所以它只以「$2$-形式的上同调类」出现在答案里；但它可能不是某个整体 $1$-形式的导数，就像磁单极的矢势写不成整体 $1$-形式。}
\pitfall{不是规范场：$B$ 没有自己的动力学，是背景几何数据，也没有非阿贝尔版本。不能总把它写成 $d\lambda$：整体上它可能是非平凡 gerbe 的联络，这正是 §7 中「扭化」的来源。它也不是 $\omega$ 的虚部：两者独立，只有 A-模型依赖复化组合 $[B+i\omega]$。}
\end{word}

\begin{word}{余各向同性}{coisotropic}{coisotropic}{§11.1 p.185}
\mathdef{辛流形 $(X,\omega)$ 的子流形 $Y$ 称余各向同性，若对每个 $y\in Y$，辛正交 $T_yY^{\perp_\omega}=\{v\in T_yX:\omega(v,\cdot)|_{T_yY}=0\}$ 包含在 $T_yY$ 中，即 $TY^{\perp_\omega}\subset TY$。等价说法：消没理想对 Poisson 括号封闭（$Y$ 由第一类约束定义）；独立 Poisson 对易函数最多 $\frac12\dim X$ 个，故 $\dim Y\ge\frac12\dim X$。等号成立时 $Y$ 是 Lagrange 子流形（$\omega|_Y=0$ 且维数一半），所以 Lagrange 是余各向同性的特殊情形。}
\mathex{取 $X=\mathbb R^4$、$\omega=dq_1\wedge dp_1+dq_2\wedge dp_2$：$Y=\{q_1=0\}$ 是实维 $3$ 的余各向同性（$TY^{\perp_\omega}=\mathrm{span}\{\partial_{p_1},\partial_{p_2}\}\subset TY$），$Y=\{q_1=q_2=0\}$ 是 Lagrange。取 $X=T^*M$、$\omega$ 为典范形式：$\pi^{-1}(S)=T^*_SM$ 对任意子流形 $S\subset M$ 都余各向同性，$S$ 为点时它是 Lagrange 纤维，$S=M$ 时它就是 $X$ 本身。}
\physdef{A-膜可以支撑在余各向同性子流形上（不限于 Lagrange）。支撑 $Y$ 上的数据是一个带联络的酉丛，其曲率 $F$ 须满足由 $Y$ 决定的方程：$Y$ 是 Lagrange 时条件是 $F+B|_Y=0$（$B$ 为 $B$-场）；$Y=X$ 时条件是 $(\omega^{-1}F)^2=-1$，即 $\omega^{-1}F$ 是近复结构（§11.1，式 11.1）。后一类膜只在 $\omega$ 是某个全纯辛形式的实部时才存在。}
\physex{在 $\mathcal M_H$ 上取 $\omega=(\mathrm{Im}\,\tau)\omega_K$：Hitchin 纤维是复维 $\dim B$、实维恰为 $\frac12\dim\mathcal M_H$ 的 Lagrange 子流形，所以也是余各向同性的特例（type F 膜就支撑在它上面）；$Y=\mathcal M_H$ 本身给出典范余各向同性膜（§11.1）。§5.5 中「固定 $E$ 让 $\phi$ 变」的一族是复 Lagrange 子流形，属于 $(B,A,A)$ 型。}
\bridge{余各向同性是「辛意义下的子对象」：它由 Poisson 对易关系（第一类约束）定义，因而在数学上对应 Poisson 几何的余各向同性子流形、在物理上对应能容纳 A-膜的支撑。消去沿 $Y$ 的零方向（$\ker\omega|_Y$ 的叶状结构）得到辛流形，$Y$ 上的 A-膜数据可视为约化后的对象，这是 §11 把 A-膜与可积系统联系起来的入口。}
\intuit{各向同性指辛形式在子流形上为零、维数不超过一半；余各向同性正好反过来：辛正交方向落在子流形内部，所以它「足够大」，大到约束可以 Poisson 对易。}
\pitfall{不是各向同性：各向同性要求 $TY\subset TY^{\perp_\omega}$、$\dim Y\le\frac12\dim X$；余各向同性要求反向包含、$\dim Y\ge\frac12\dim X$，两者只在 Lagrange 处重合。也不是随便一个 Poisson 交换函数的零点集：必须取到最大数量的独立第一类约束，否则维数会掉到中维以下。}
\end{word}

\begin{word}{A-膜}{A-brane}{a-brane}{§5.2 p.69}
\mathdef{辛流形 $(X,\omega)$（带 $B$-场）上 A-模型的对象：支撑 $Y\subset X$ 余各向同性，Chan-Paton 为一个带联络的酉丛，曲率满足边界条件给出的方程。当 $Y=L$ 是 Lagrange 子流形时，条件是 $\dim L=\frac12\dim X$、$\omega|_L=0$，且 $L$ 上的 $U(1)$ 丛平坦（$F+B|_L=0$）；这样的 $(L,\nabla)$ 就是通常说的拉格朗日膜，也就是 Fukaya 范畴的对象。态空间由两端落在两个膜上的开弦态给出，数学上是 Floer 型上同调。}
\mathex{取 $X=T^2$、$\omega=\mathrm{vol}$：$L=S^1$ 是 Lagrange 子流形，加平坦 $U(1)$ 联络 $\nabla=d+i\theta\,dt$ 给出一族 A-膜，参数 $\theta\in\mathbb R/\mathbb Z$。取 $X=\mathbb C^2$、$\omega$ 为标准形式：$L=\mathbb R^2=\{z_1,z_2\in\mathbb R\}$ 是 Lagrange 子流形，平凡丛给 A-膜；$X=T^*M$ 的零截面 $M$ 也是 Lagrange 的，其上 A-膜在镜像侧对应 $T^*M$ 上的相干层。}
\physdef{A-膜是世界面 $\Sigma$ 上带边理论的边界条件，且要求边界条件保持 A-模型的拓扑超荷 $Q$。世界面上它给出类似 $G_{IJ}\partial_1\Phi^J=-iF_{IJ}\partial_0\Phi^J$ 的边界方程（式 11.6），其经典解由 $\omega$ 与 $F$ 的非退化性强烈约束（§11.1）。在 KW06 中 A-膜不一定是 Lagrange 的，$Y=X$ 的典范余各向同性膜就是一例。}
\physex{在 $\mathcal M_H$ 上，Hitchin 纤维 $F_b$ 是 Lagrange 子流形（关于 $\omega_J$、$\omega_K$ 都是），加平坦 $U(1)$ 丛就是 type F 膜，属于 $(B,A,A)$ 型：它在 $I$ 中全纯，关于 $\omega_J$、$\omega_K$ Lagrangian；它是零膜在 $S$-对偶下的像，也是 §9 中磁特征膜的原型（\S5.4--\S5.5）。}
\bridge{A-膜对应辛几何对象；B-膜对应复几何对象（复子流形加全纯丛）。同调镜像对称把 $X$ 的 A-膜范畴与镜像 $\hat X$ 的 B-膜范畴等价起来。在 KW06 的具体框架里还有第三步：用典范余各向同性膜 $A_{cc}$ 作函子 $A\mapsto\mathrm{Hom}(A_{cc},A)$，把 A-膜送到 $G$-丛模空间上的 D-模（\S11.1 与 §9）。}
\intuit{A 是「面积」：A-膜像一张只感知辛面积的膜，它的位置由 Lagrange 或余各向同性的条件决定，$\omega$ 在它上面的限制尽量小。靠得住是因为 A-模型只依赖 $\omega$ 与 $B$ 的组合。}
\pitfall{A-膜不是 A-模型：前者是对靶空间的边界条件（范畴的对象），后者是世界面上的二维场论。也不要求支撑一定是 Lagrange 子流形：余各向同性的支撑也可以，这正是 §11 处理的情形；但支撑不能是任意的复子流形。}
\end{word}

\begin{word}{B-膜}{B-brane}{b-brane}{§5.3 p.70}
\mathdef{复流形 $(X,J)$ 上 B-模型的对象：支撑 $N\subset X$ 是复子流形，Chan-Paton 是 $N$ 上的全纯向量丛或更一般的相干层 $\mathcal E$；作为 $X$ 上的对象，它就是推前层 $i_*\mathcal E$。态空间是 $\mathrm{Ext}^\bullet_{\mathcal O_X}(i_*\mathcal E,i_*\mathcal F)$，即全纯态射的上同调；所有 B-膜连同态射构成 $X$ 上相干层的导出范畴 $D^b(\mathrm{Coh}(X))$。}
\mathex{取 $X=T^*M$、复结构为标准复结构：支撑零截面 $M$ 的 B-膜、Chan-Paton 为 $M$ 上的全纯向量丛 $\mathcal E$，整体对应 $M$ 上的相干层；反过来，$T^*M$ 上任何支撑在零截面的相干层都拉回到 $M$。取 $X=\mathbb P^1$：填满空间的膜带 $\mathcal O(k)$；支撑在点 $p$ 的膜带天穹层 $\mathcal O_p$，且 $\mathrm{Ext}^1(\mathcal O_p,\mathcal O_p)=\mathbb C$ 对应点沿 $\mathbb P^1$ 移动。}
\physdef{B-膜是 B-模型的边界条件：要求 $\Phi$ 把世界面的边界映进复子流形 $N$，且 Chan-Paton 丛限制到 $N$ 上全纯（$\bar\partial_E=0$）。$J_+=J_-=\hat J$ 时得到复结构 $\hat J$ 中的 B-模型（$t=\pm i$），其膜就是上述全纯对象（\S5.1、\S5.3）。$B$-场的 $(0,2)$ 分量把全纯丛扭化成扭化层，膜的范畴随之改变（§7）。}
\physex{在 $\mathcal M_H$ 的复结构 $J$ 中，B-膜是 $J$-全纯的对象；但 Hitchin 纤维在 $J$ 中并非复子流形（它关于 $\omega_J$ 是 Lagrange 的），所以 type F 膜不是 $J$ 中的 B-膜，而是 $I$ 中的 B-膜——这正是「(B,A,A) 型」的第一字母指 $I$ 的意思（\S5.5）。$(B,B,B)$ 型的例子是零膜与填满空间带平凡丛的膜。}
\bridge{对象级：相干层 $\mathcal F$ 映到支撑在 $\mathrm{supp}\,\mathcal F$ 上、Chan-Paton 为 $\mathcal F$ 的膜。态射级：两端落在 $\mathcal E,\mathcal F$ 上的开弦态空间等于 $\mathrm{Ext}^\bullet(\mathcal E,\mathcal F)$，鬼数即上同调次数；开弦算子的乘积展开对应 Ext 的合成。导出范畴中的拟同构就是物理上被 BRST 零模除掉的态。Fourier--Mukai 变换由 $X\times\hat X$ 上的核层给出（§5.1、§9）。}
\intuit{B-膜是一块全纯的膜：只在意复结构、不在意度量，所以能在全纯世界里变形；丛一旦忘了全纯结构就不再有意义。}
\pitfall{B-膜不是 B-模型：前者是边界条件（范畴对象），后者是世界面的二维场论。也不等于「任意实子流形加丛」：必须是复子流形加全纯丛。B-膜一般不是 A-膜：支撑在 Lagrange 子流形上的膜通常是 A-膜，除非它同时是复子流形且维数恰为一半。}
\end{word}

\begin{word}{膜}{brane}{brane}{§1 p.2}
\mathdef{带边世界面 $\Sigma$ 的边界条件对应的靶空间对象：数据为支撑 $Y\subset X$ 与 Chan-Paton 数据（$Y$ 上的丛、层或层复形，带联络或全纯结构），并保持拓扑超荷。数学上膜就是范畴的对象：$A$-模型的膜是余各向同性子流形加丛，$B$-模型的膜是复子流形加全纯层；两膜之间的态射是开弦态空间，膜的模空间即形变空间。}
\mathex{最基本的三个例子：填满靶空间的膜（$Y=X$，Chan-Paton 为平凡线丛），它在范畴中是单位对象；支撑在一点 $p$ 的零膜（Chan-Paton 为天穹层 $\mathcal O_p$），它是「点对象」；$T^*M$ 上支撑零截面的膜带 $M$ 上相干层，它把 $M$ 上的层论搬到 $T^*M$。它们分别对应单位、点与推前函子。}
\physdef{膜是带边二维理论的边界条件；通常由 D-膜实现，膜上住着 Chan-Paton 丛的联络描述的规范场。KW06 约定：不特指时写「膜」而不写「D-膜」，因为扭化后的理论里膜已不是通常的 D-膜，而支撑与 Chan-Paton 丛才是它的两个基本数据（\S1、\S5.5）。}
\physex{KW06 §5.5 的三类具体膜：支撑在点 $p$ 的零膜 $B_p$；支撑整个 $\mathcal M_H$、带平凡 Chan-Paton 丛的膜；支撑 Hitchin 纤维 $F_b$、带平坦 $U(1)$ 丛的 type F 膜。它们分别属于 $(B,B,B)$、$(B,B,B)$、$(B,A,A)$ 型；$S$-对偶作用在膜上，把 $(B,B,B)$ 型换成 $(B,A,A)$ 型（\S5.4）。}
\bridge{膜对应范畴的对象。数学侧：膜对应相干层（$B$-模型）或 Lagrange、余各向同性对象（$A$-模型）；膜之间的开弦态空间对应 $\mathrm{Ext}$ 群或 Floer 型上同调；膜的形变对应 $\mathrm{Ext}^1$。罗塞塔词典的骨架是一条范畴等价，膜正是两边对象。}
\intuit{膜是「世界面可以停下来的地方」：边界落在支撑上，还带着内部自由度，像弦的端点能在膜上滑动并携带电荷。}
\pitfall{不是 D-膜的同义词：原文的膜可以带非平凡 $B$-场扭化、可以支撑在余各向同性子流形上，这些不是通常 D-膜的语言。也不能把「膜」与「模型」混为一谈：膜是对靶空间的边界条件，模型是世界面上的场论。}
\end{word}

\begin{word}{边界条件}{boundary condition}{boundary-condition}{§5.5 p.74}
\mathdef{世界面 $\Sigma$ 的边界上的局部条件：指定 $\Phi(\partial\Sigma)$ 所在的子流形 $Y\subset X$，并给出法向导数与切向导数的关系（如 $G_{IJ}\partial_1\Phi^J=-iF_{IJ}\partial_0\Phi^J$，$F$ 为 Chan-Paton 联络的曲率），使边界项与体作用量一起保持超对称。能作为「膜」的边界条件必须保持拓扑超荷 $Q$；保持不同超荷的选择给出不同的膜类型。}
\mathex{取 $\Sigma$ 为上半平面、$X=T^2$、$Y=S^1\times\{0\}$：边界条件要求 $\Phi(\partial\Sigma)\subset S^1$，切向导数的自由度由平坦 $U(1)$ 联络的 holonomy $\theta$ 决定，这给出 $T^2$ 上一族 A-膜。取 $X=\mathbb C$ 的 B-模型、$Y=\{0\}$：要求 $\Phi$ 在边界附近全纯且落在原点，给出零膜。$T^*M$ 上 Neumann 条件（$\partial_1\Phi=0$ 型）对应的膜与 $M$ 上相干层相关。}
\physdef{边界条件局部由支撑与 Chan-Paton 丛描述，整体才给出一个膜（\S5.5）。拓扑不变量只看保持 $Q$ 的边界条件；保持一个超荷给相应模型的一个膜，保持更多超荷则它在多个复结构中同时是 A 型或 B 型，这就是膜的类型。}
\physex{在 $\mathcal M_H$ 上，$(B,B,B)$ 型的膜（零膜、填满空间的膜）与两个线性无关的拓扑超荷组合相容；$(B,A,A)$ 型（type F 膜）与另一组相容，其组合分别生成 $I$ 中 B-模型与 $J$、$K$ 中 A-模型的超对称。$S$-对偶交换这两组，从而交换膜类型（\S5.4）。}
\bridge{边界条件等同于膜，等同于范畴的对象；两个边界条件之间的开弦态就是范畴的 $\mathrm{Hom}$。模型的选择决定边界条件里哪些几何数据进入：A-模型用 $(Y,\omega,B)$，B-模型用 $(N,J,\text{全纯丛})$。所以同一子流形在不同模型里可以是不同的膜。}
\intuit{边界条件规定了弦的端点能去哪里、端点能带什么电荷、法向能不能动；用物理的话说，它就是膜对弦的「邀请函」。}
\pitfall{不是任意局部条件都能给膜：必须保持拓扑超荷，否则只是普通边界条件，不给出拓扑不变量。原文也不把「局部边界条件」与「整体膜」当同一层次：膜是整体对象，边界条件只是局部描述。}
\end{word}

\begin{word}{Chan-Paton 丛}{Chan-Paton bundle}{chan-paton-bundle}{§6.4 p.95}
\mathdef{膜上的向量丛（更一般地：层或层复形），通常带联络 $\alpha$；其纤维是弦端点携带的规范量子数所在的空间。数学上在 B-模型侧它就是膜对应的相干层，在 A-模型侧是带酉联络的丛。把膜与固定向量空间 $V$ 张量积，就是把 Chan-Paton 丛换成 $E\otimes V$，这对膜的其他数据没有影响。}
\mathex{$T^2$ 上的 A-膜（$L=S^1$）：Chan-Paton 是 $\mathbb C$ 上的平坦 $U(1)$ 线丛，参数为 holonomy $\theta\in\mathbb R/\mathbb Z$。$\mathbb P^1$ 上的 B-膜：全纯线丛 $\mathcal O(k)$ 是秩 $1$ 的 Chan-Paton 层；零膜：Chan-Paton 是支撑在点的天穹层，秩为 $1$ 但纤维只在点上非零。$T^*M$ 上支撑零截面的膜：Chan-Paton 为 $M$ 上秩 $r$ 的全纯向量丛 $E$。}
\physdef{膜上住规范场，Chan-Paton 丛负责「颜色」：规范群作用在它的纤维上。Wilson 线作用在膜上就是把 Chan-Paton 丛与泛丛 $E_q(R)$ 张量积，$U\mapsto U\otimes E_q(R)$（§8）；特征膜的条件是 $U\otimes E_q(R)\cong U\otimes V$（$k$ 个拷贝），其中 $V$ 是固定向量空间，所以「特征值」应想成向量空间 $V$ 而不是维数 $k$（§6.4）。}
\physex{KW06 §6.4（p.95）用 $k$ 个 Chan-Paton 拷贝与 $V$ 的张量积说明这一点；零膜的 Chan-Paton 层是支撑在 $x$ 的天穹层 $U_x$，限制到点后一切丛都平凡，故特征值只是固定向量空间 $E_q(R)|_x$。在 §7 中，平坦 $B$-场 $b\in H^2(\mathcal M_H,U(1))$ 使 Chan-Paton 丛变成扭化向量丛（扭化层）。}
\bridge{Chan-Paton 丛就是膜的「层」那一半数据：数学侧的对象是（支撑，层），其不变量是层上同调；A-模型侧换成带联络的酉丛。平坦 $B$-场对应扭化层，这是 B-场具体进入膜的世界的方式（§7、§11.1）；在导出范畴语言里，零膜的 Chan-Paton 是天穹层 $\mathcal O_x$ 而不是零维向量空间。}
\intuit{弦的端点像电荷：颜色指标住在 Chan-Paton 丛的纤维里，端点移动时这个向量空间跟着平移，联络就是平移的规则。}
\pitfall{不是背景规范场：Chan-Paton 丛是膜的内部数据，规范场只活在膜上。也不要只想成向量丛：原文明确允许更一般的层或层复形（导出范畴的对象），零膜的 Chan-Paton 就是天穹层。扭化情形下它甚至不是通常的丛，而是扭化层。}
\end{word}

\begin{word}{零膜}{zerobrane}{zerobrane}{§5.5 p.74}
\mathdef{支撑在一点 $p\in X$、Chan-Paton 为该点天穹层 $\mathcal O_p$ 的膜。在 $B$-模型中它就是点对象 $\mathcal O_p\in D^b(\mathrm{Coh}(X))$，自同态代数为 $\mathrm{Ext}^\bullet(\mathcal O_p,\mathcal O_p)=\wedge^\bullet T_pX$，一阶形变空间就是切空间 $T_pX$。它是「最小的膜」，但并不平凡：点的结构层把靶空间该点的切方向都记了下来，因此它是范畴里的点对象、也是电特征膜的原型。}
\mathex{取 $X=\mathbb P^1$：零膜 $\mathcal O_p$ 满足 $\mathrm{Ext}^1(\mathcal O_p,\mathcal O_p)=\mathbb C$，几何意义是点可以沿 $\mathbb P^1$ 移动。取 $X=T^2$：$\mathrm{Ext}^1=\mathbb C^2$，对应点在二维环面上的两个移动方向。取 $X=\mathcal M_H$：零膜 $B_p$ 支撑在光滑点 $p$，它的模空间是 Hitchin 纤维 $F_b$（\S5.4 的二维解释），这一句是 §5.4 的关键。}
\physdef{零膜作为边界条件：世界面的边界收缩到一点 $p$，$\Phi(\partial\Sigma)=p$，Chan-Paton 只有该点的天穹层。它是 $(B,B,B)$ 型——点在任意复结构里都是（零维）复子流形，所以它在 $I,J,K$ 三个复结构中都是 B-膜，不是 A-膜。$S$-对偶后它变成 $(B,A,A)$ 型的膜，支撑在对偶 Hitchin 纤维上（\S5.4--\S5.5）。}
\physex{在 $\mathcal M_H$ 上，零膜在任意 $e_0$ 都存在（§7）；它是 Wilson 线算子的电特征膜（§8），其模空间就是 Hitchin 纤维 $F_b$（\S5.4）；$S$-对偶后对偶膜的模空间是 $J_{\tilde F_b}$，维数相同。物理上它是「电」这一侧最简单的对象，$S$-对偶把它送到「磁」侧的 type F 膜。}
\bridge{零膜对应点层 $\mathcal O_p$（$B$-模型侧）；对象级与态射级都吻合：$\mathrm{Ext}^\bullet(\mathcal O_p,\mathcal O_p)=\wedge^\bullet T_pX$ 就是驻在 $p$ 处的开弦态空间，形变即切向量。在几何 Langlands 的版本里，平坦 ${}^LG$-丛 $E$ 的天穹层 $\mathcal O_E$ 映到 $A$-膜 $(F_b,\nabla_E)$，即 Hitchin 纤维加对应的平坦 $U(1)$ 丛——这是「零膜映到辛对象」的精确陈述。}
\intuit{零膜像一根钉子把弦的端点钉在靶空间的一点上；钉子虽小，但钉住的点的切空间成了它的内部自由度，所以它不是「无」。}
\pitfall{不是「什么都不带」的平凡对象：自同态是外代数而不是 $\mathbb C$，形变是切空间。也不是 A-膜：它在三个复结构下都是 B 型，A-侧的对应物要经 $S$-对偶或镜像才能得到（那时它才变成辛意义下的对象）。}
\end{word}

\begin{word}{type F 膜}{type F brane}{type-f-brane}{§9 p.125}
\mathdef{支撑在 Hitchin 纤维 $F_b\subset\mathcal M_H$ 上、带秩 $1$ 平坦 $U(1)$ Chan-Paton 丛的膜。$F_b$ 是 Hitchin 纤维化 $\pi:\mathcal M_H\to B$ 的纤维，在 $I$ 中是复阿贝尔簇（torsor），关于 $\Omega_I$ 复 Lagrangian，从而关于 $\omega_J$、$\omega_K$ 也是 Lagrange 子流形；所以它是 $(B,A,A)$ 型：在 $I$ 中全纯，关于 $J$、$K$ 是 A-膜。数学上它由「基上的点 $b$ 加纤维上平坦 $U(1)$ 丛」给出，模空间是 $B$ 上配一束对偶环面。}
\mathex{取 $G=SU(2)$、$C$ 亏格 $2$：$B=H^0(C,K_C^2)$ 复维 $3$，$\mathcal M_H$ 复维 $6$，纤维 $F_b$ 是复维 $3$ 的阿贝尔簇（谱曲线的广义 Prym），关于 $\omega_J$、$\omega_K$ Lagrange。$F_b$ 上的平坦 $U(1)$ 丛由 $J_{F_b}=H^1(F_b,\mathbb R)/H^1(F_b,\mathbb Z)$ 参数化，type F 膜的模空间就是这个对偶环面丛。}
\physdef{type F 膜是零膜的 $S$-对偶，因而预期是磁特征膜，即 't Hooft 线算子的特征膜；它是几何 Langlands 观点下最重要的一类 A-膜（§9）。「F」指 Hitchin 纤维（fiber），不是别的；它保持的拓扑超荷组合给出 $I$ 中的 B-模型与 $J$、$K$ 中的 A-模型超对称。}
\physex{\S5.4 的二维解释：零膜 $B_p$ 是 $(B,B,B)$ 型；$S$-对偶后对偶膜支撑在对偶 Hitchin 纤维 $\tilde F_b$ 上，是 $(B,A,A)$ 型。又因 A-膜的支撑维数至少是中间维、而 Hitchin 纤维维数至多是中间维，故支撑恰为 $\tilde F_b$——这正是 type F 膜。零膜的模空间 $F_b$ 与对偶膜的模空间 $J_{\tilde F_b}$ 维数相同。}
\bridge{type F 膜对应「Hitchin 纤维加平坦 $U(1)$ 丛」，也就是谱曲线/Prym 簇上的一个点；在几何 Langlands 中它对应 Hecke 特征层（$D$-模）在膜一侧的像，$S$-对偶把 Wilson 函子与 Hecke 函子互换（§8--§9）。所以 type F 膜是「膜」这一列与「$D$-模」这一列接口处的对象。}
\intuit{type F 膜是缠在 Hitchin 纤维这个环面上的膜：它沿纤维铺开并带一个平坦相位，看不见纤维内部的复结构、只看得到环面结构，所以在 $J$、$K$ 中是辛/A 型的。}
\pitfall{不是某一单个 Lagrange 膜：它是整个纤维加上平坦 $U(1)$ 丛的一整族，模空间是环面丛。也不是 $(A,A,A)$ 型（不存在）：它的第一字母是 B，在 $I$ 中全纯这一点最易漏。它更不是「填满空间的膜」，虽然两者在 $S$-对偶下配对。}
\end{word}

\begin{word}{典范余各向同性 A-膜}{canonical coisotropic A-brane}{canonical-coisotropic-a-brane}{§11.1 p.185；§1 p.5}
\mathdef{取 $X=\mathcal M_H$、$\omega=(\mathrm{Im}\,\tau)\omega_K$、支撑 $Y=X$、Chan-Paton 为平凡线丛 $N$ 带联络 $\alpha$，要求曲率满足 $F=(\mathrm{Im}\,\tau)\omega_J$（式 11.2）。此时 $(\omega^{-1}F)^2=-1$（式 11.1）恰好成立：$N=\omega_K^{-1}\omega_J=I$ 是复结构。由 $\omega_J$ 恰当，$F$ 确实是平凡线丛上联络的曲率；这样得到的 A-膜称典范余各向同性膜 $B_{cc}$。$G$ 单连通时唯一，否则唯一到 Chan-Paton 丛扭化一个平坦线丛。}
\mathex{更一般的解是 $F=(\mathrm{Im}\,\tau)(a\omega_J+b\omega_I)$，$a^2+b^2=1$（式 11.4）；但 $b\ne 0$ 时 $\omega_I$ 不恰当、周期不是 $2\pi$ 的整数倍，$F$ 一般不是任何联络的曲率，因此只有 $a=\pm1$ 的两个解可用。在 $\mathcal M_H$ 含 $T^*M_{st}$ 为 Zariski 开集的意义下，$B_{cc}$ 与 $T^*M_{st}$ 上的特殊线丛限制后，其自同态代数给出 $M_{st}$ 上的微分算子层（§11.1、§1 注）。}
\physdef{c.c. 膜是 $Y=X$ 的余各向同性 A-膜：支撑整个 $\mathcal M_H$，Chan-Paton 是 $U(1)$ 线丛，条件 $(\omega^{-1}F)^2=-1$。它是 $(A,B,A)$ 型：在 $I$、$K$ 中是 A-膜，在 $J$ 中是 B-膜；验证用 $(\omega_I^{-1}\omega_J)^2=-1$（式 11.5）以及「曲率是 $\omega_J$ 的倍数，而 $\omega_J$ 在 $J$ 中是 $(1,1)$ 型」。}
\physex{经典边界算子层是 $\mathcal M_H$ 上关于 $I$ 的全纯函数层（鬼数 $0$ 的可观测量），局部无高阶上同调、不可形变、无量子修正，整体可能有量子修正。存在 $F$ 时开弦代数非交换，参数是 $(g+F)^{-1}$ 的反对称部分，与形变量子化直接相关。世界面瞬子情形：边界条件 $G_{IJ}\partial_1\Phi^J=-iF_{IJ}\partial_0\Phi^J$（式 11.6）与全纯方程合起来给 $(\omega-iF)\partial_0\Phi=0$（式 11.7），$\omega$、$F$ 非退化，故边界上的瞬子只能是常数，无瞬子修正。}
\bridge{c.c. 膜是把 A-膜翻译成 $D$-模的机器：函子 $A\mapsto\mathrm{Hom}(A_{cc},A)$ 把 $\mathcal M_H$ 上的 A-膜送到 $M_{st}(G,C)$（$G$-丛模空间）上的 $D$-模，理由是 $\mathrm{Hom}(A_{cc},A_{cc})$ 给出微分算子层。这正是 Witten 2008 与 Frenkel 2009 所述的链：B-膜（${}^LG$ 侧）$\leftrightarrow$ A-膜（$\mathcal M_H$ 侧）$\leftrightarrow$ $D$-模（$\mathrm{Bun}_G$ 侧）。}
\intuit{c.c. 膜像一整块盖住模空间的膜：曲率与 $\omega_J$ 成比例，恰好造出与 $I$ 相同的复结构，故它既满足辛条件、又带全纯结构，才能把两侧对象接起来。}
\pitfall{不是某个具体的 Lagrange 膜，也不是「所有膜都与它同构」：它的支撑是全体 $\mathcal M_H$。这类膜只在 $\omega$ 是某全纯辛形式实部时存在，且公式只在 $\mathcal M_H$ 光滑部分成立（整体由底层规范理论的边界条件定义）。它也不是 $(A,A,A)$ 型（不存在）：在 $J$ 中它是 B-膜。}
\end{word}

\begin{word}{膜的类型}{brane type}{brane-type}{§5.5 p.75}
\mathdef{对 $\mathcal M_H$ 的三个复结构 $I,J,K$ 分别标注该膜是 A 型（辛意义：支撑余各向同性、带联络丛）还是 B 型（全纯意义：支撑复子流形、带全纯层），得到三字母串 $(x,y,z)$，依次对应 $I,J,K$。一个膜能同时在多个复结构下是 A 型或 B 型，因为同一个支撑与同一个丛可以同时满足不同结构的条件；类型就是这种「多重身份」的记账方式。}
\mathex{四种可能全部出现：$(B,B,B)$，例为零膜与填满空间带平凡丛的膜；$(B,A,A)$，例为 Hitchin 纤维带平坦 $U(1)$ 丛（关于 $\Omega_I$ 复 Lagrangian）；$(A,B,A)$ 与 $(A,A,B)$，例为固定 $A_{\bar z}+i\phi_{\bar z}$ 或 $A_{\bar z}-\phi_{\bar z}$ 得到的复 Lagrange 族。不存在 $(B,B,A)$：$I$、$J$ 都全纯意味着 $K=IJ$ 也全纯；$(A,A,A)$ 同样不存在。}
\physdef{膜的类型由它保持哪些拓扑超荷决定：$(B,B,B)$ 型保持两个线性无关的拓扑超荷，其组合生成 $I$、$J$、$K$ 中的 B-模型；$(B,A,A)$ 型保持另一组，其组合生成 $I$ 中 B-模型与 $J$、$K$ 中 A-模型的拓扑超对称。$S$-对偶交换 $(B,B,B)\leftrightarrow(B,A,A)$，保持 $(A,B,A)$ 与 $(A,A,B)$ 不变（\S5.4--\S5.5，Table~2）。}
\physex{在 $\mathcal M_H$ 上：零膜 $B_p$ 与填满空间的膜是 $(B,B,B)$；type F 膜与「固定 $E$ 让 $\phi$ 变」的族是 $(B,A,A)$；固定 $A_{\bar z}\pm i\phi_{\bar z}$ 的族分别给 $(A,B,A)$ 与 $(A,A,B)$。$S$-对偶把 $(B,B,B)$ 的零膜送到 $(B,A,A)$ 的 type F 膜，这正是四维 $S$-对偶在二维膜上的表现。}
\bridge{三个字母对应三个复结构 $I,J,K$，不是三种膜、也不是三个模型；A/B 分别指「用辛形式定义」与「用复结构定义」该膜。这个三元组把「同一个膜在不同复结构下的身份」压成三个字母，于是 $S$-对偶、镜像、$\varphi\mapsto i\varphi$ 等操作对膜的作用可以一眼算出（\S5.4--\S5.5）。}
\intuit{像给同一块膜拍三张 X 光片：每张检查它在某个复结构下是全纯（B）还是辛（A），三张叠起来就是它的类型。}
\pitfall{不是 A-模型/B-模型的标签，也不是「膜有三种类型」：字母对应三个复结构，一个膜可以同时是 A 型和 B 型（例如 $(B,A,A)$）。也不存在 $(B,B,A)$ 与 $(A,A,A)$ 两种看似自然的组合。}
\end{word}

\begin{word}{SYZ 图景}{SYZ picture}{syz-picture}{§5.5 p.77}
\mathdef{Strominger--Yau--Zaslow 提出的镜像对称几何机制：一对镜像 Calabi--Yau 流形 $X$ 与 $\hat X$ 各带特殊 Lagrange 环面纤维化 $\pi:X\to B$、$\hat\pi:\hat X\to B$，同一基上的纤维 $F_b$ 与 $\hat F_b$ 互为对偶环面（$H_1(F_b,\mathbb Z)$ 与 $H^1(F_b,\mathbb Z)$ 对偶），且 $X$ 的复结构模与 $\hat X$ 的复化 Kähler 模互换。镜像对称被解释为「基上同一映射加纤维上的 $T$-对偶」。}
\mathex{取 $X=E_\tau=\mathbb C/(\mathbb Z+\tau\mathbb Z)$：纤维化可取纤维 $S^1$、基 $S^1$，纤维上的 $T$-对偶把 $\tau\mapsto-1/\tau$，于是复结构模与复化 Kähler 模互换。取 $\mathcal M_H$ 情形：Hitchin 纤维化 $\pi:\mathcal M_H\to B$ 的纤维是复阿贝尔簇（torsor），对偶纤维化是 ${}^LG$ 的 Hitchin 纤维化，基可自然等同（单连通型时 $\Xi=\mathrm{id}$），纤维上的 $T$-对偶实现镜像对称。}
\physdef{SYZ 图景解释四维 $S$-对偶在二维的像：$S$-对偶在 Hitchin 基 $B$ 上作用且保分次环，在纤维上无法用局部算子实现，只能是 $T$-对偶；于是零膜（$(B,B,B)$ 型）在 $S$-对偶后成为支撑对偶纤维的 $(B,A,A)$ 型膜，即 type F 膜（\S5.4--\S5.5）。}
\physex{取 $G=SU(2)$、对偶群 $SO(3)$、$C$ 亏格 $2$：两个 Hitchin 纤维化同底 $B$，对应纤维互为对偶环面（广义 Prym 簇的对偶），SYZ 图景就是「基映射 $\Xi$ 加纤维 $T$-对偶」。原文引 Hausel--Thaddeus 对酉群、Hitchin 对 $G_2$、Donagi--Pantev 对一般半单群，作为纤维对偶成立的依据。}
\bridge{SYZ 图景对应镜像对称的几何实现，也对应 Fourier--Mukai 变换：纤维上的 $T$-对偶在 $B$-膜范畴里表现为核层给出的 FM 反变函子，在 $A$-膜一侧则是纤维上平坦 $U(1)$ 丛的平移。KW06 用它把四维的 $S$-对偶降成二维的 $T$-对偶（\S5.4--\S5.5）。}
\intuit{SYZ 像傅里叶变换：位置空间与动量空间互为对偶，环面纤维与对偶环面纤维就像一对傅里叶共轭的坐标；镜像对称只是把「在哪里」换成「绕多少」。}
\pitfall{原文只把它当启发式的几何图景，不当证明：一般 Calabi--Yau 上的 SYZ 猜想仍未被证明。KW06 的情形之所以可靠，是因为 Hitchin 纤维化有定理保证纤维对偶，不要把这个特例的严格性误当成 SYZ 的一般证明。}
\end{word}
